% ============================================================
% EthicalGuard v2 — LaTeX 全文（解决 G10b 三🔴：T1 定理/建模选择/权重下限）
% 编译：pdflatex/xelatex；中文翻译：EthicalGuard-paper-v2.zh.md
% 版本说明：数字以 runs 真实结果为准（KKT 1.69e-7 级、满意度 0.82-0.99）；
%          实验 4.6-4.8 为"待跑设计"（用户执行后填表），表中 TBD 占位
% ============================================================
\documentclass[11pt]{article}
\usepackage[utf8]{inputenc}
\usepackage{amsmath,amssymb}
\usepackage{graphicx}
\usepackage{booktabs}
\usepackage{multirow}
\usepackage{microtype}
\usepackage[hidelinks]{hyperref}
\usepackage{enumitem}
\usepackage{algorithm}
\usepackage{algorithmic}
\usepackage{amsthm}
\usepackage{thm-restate}

\newtheorem{theorem}{Theorem}
\newtheorem{proposition}[theorem]{Proposition}
\newtheorem{lemma}[theorem]{Lemma}
\newtheorem{corollary}[theorem]{Corollary}
\newtheorem{definition}[theorem]{Definition}
\newtheorem{remark}[theorem]{Remark}

\title{EthicalGuard: Multi-Agent Negotiation with Generalized Nash Equilibrium
for Auditable Medical Ethics Decision-Making}
\author{<作者>}
\date{\today}

\begin{document}
\maketitle

% ---------- 摘要（任务→gap→方法→结果→意义）----------
\begin{abstract}
Medical ethics dilemmas involve multiple stakeholders---physicians, patients,
families, ethics committees---whose values (autonomy, beneficence,
non-maleficence, justice) conflict and must be reconciled under resource and
safety constraints. Single-model ethical reasoning is opaque, lacks
multi-stakeholder perspectives, and provides no formal guarantee that an
ethical floor is respected. We propose \textbf{EthicalGuard}, a multi-agent
negotiation framework that formulates medical ethics decision-making as a
\textbf{generalized Nash equilibrium (GNE)} problem with a
\textbf{regularized primal--dual solver}: each stakeholder agent holds a
strategy over treatment intensity and four-principle weights, and all agents
jointly carry a shared set of coupling constraints---resource feasibility and
four-principle floors. We make the semantic claim that this shared-constraint
structure is the mathematical form of an \emph{ethical floor}: negotiation is
free within the floors; the floors are non-negotiable. We prove a
\textbf{floor-guarantee theorem} (any converged solution satisfies the floors
up to numerical tolerance, and active constraints bind at equality), and we
justify the choice of GNE over Nash bargaining, Stackelberg, and cooperative
games. To prevent low-expressiveness stakeholders from being silenced, we
introduce a \textbf{reliability-weight floor} that decouples information
reliability from stance representativeness. On five medical-ethics datasets,
negotiation converges with KKT residuals on the order of $10^{-7}$ and mean
stakeholder satisfaction above 0.8 (0.82--0.99 across scenarios); ablations
confirm each component's contribution; stress tests show recovery under
pressure and document when principles are abandoned; and floor-activation,
collusion, and fairness experiments are designed to validate the theoretical
guarantees.
\end{abstract}

\begin{keywords}
medical ethics \sep multi-agent negotiation \sep generalized Nash equilibrium
\sep shared constraints \sep ethical floor \sep fairness
\end{keywords}

% ---------- 1. 引言（骨架 S1）----------
\section{Introduction}
% P1 背景漏斗
Medical large language models (LLMs) are increasingly deployed in clinical
decision support, yet ethical dilemmas---when a patient's autonomy conflicts
with a physician's beneficence, or a family's wishes strain scarce
resources---cannot be reduced to a single model's next-token prediction
\cite{challenge2025,frontiers2026}. Such dilemmas are inherently
multi-stakeholder: different parties hold different values (autonomy,
beneficence, non-maleficence, justice), and the decision must be reconciled
under hard resource and safety constraints. A single model reasoning about
ethics is opaque, has no mechanism to represent genuinely conflicting
perspectives, and---most dangerously---offers no formal guarantee that a
minimal ethical standard (an \emph{ethical floor}) is respected.

% P2 缺口制造 + 实例
Prior work addresses pieces of this problem: uncertainty-aware diagnostic
models quantify ``not knowing'' \cite{confidx,entropy2026}; cost-sensitive
decision frameworks optimize over money and harm \cite{graphdx}; and
multi-agent LLM systems assemble heterogeneous perspectives
\cite{aegle,emomas,kamac,graphofstates}. Most directly,
\emph{Many LLMs Are More Utilitarian Than One} \cite{utilitarian2025} shows
that aggregating multiple LLMs shifts group decisions toward utilitarianism.
Yet aggregation and voting neither model \emph{strategic} stakeholders nor
enforce \emph{shared constraints}: a majority can still converge to a
floor-violating outcome. Contrastively, when shifting to the perspective of
\emph{negotiation under a non-negotiable floor}, a critical demand emerges,
which we call the \textbf{ethical floor as a shared constraint}: the floor
must be part of the equilibrium problem itself---negotiated around, but never
crossed. As shown in Fig.~\ref{fig:motive}, consider the detained trauma
patient BB whose autonomy (privacy from a correctional officer) conflicts with
the physician's urgent beneficence; five stakeholders must reach an agreement
that respects the autonomy floor while enabling life-saving care.
\begin{figure}[t]
  \centering
  % \includegraphics[width=\columnwidth]{motive.png}
  \caption{Motivation: the detained trauma patient BB---a single-model hard
  decision ignores the autonomy floor, whereas multi-stakeholder negotiation
  must reconcile beneficence with the non-negotiable autonomy floor.}
  \label{fig:motive}
\end{figure}

% P3-P4 挑战分解
We empirically conclude two challenges when formalizing this demand:
\begin{enumerate}[label=(\Roman*)]
  \item \textbf{The floor must be provably enforced, not merely
  asserted}：prior ``constitution''-style approaches \cite{superego2025,
  governance2025} apply normative constraints statically, outside the
  negotiation dynamics; a static filter offers no equilibrium-level guarantee.
  We need a \emph{theorem} that any converged outcome satisfies the floors,
  and evidence that the constraint is \emph{active} (not trivially satisfied).
  \item \textbf{The choice of game-theoretic solution concept must be
  justified, not assumed}：Nash bargaining has axiomatic foundations, yet
  cannot express shared non-negotiable floors; Stackelberg presupposes an
  authority hierarchy that is itself contested; cooperative games assume
  transferable utility that principle weights do not satisfy. Without an
  explicit modeling-choice argument, ``ethics as GNE'' remains a metaphor.
  In addition, a mechanism that weights stakeholders by reliability can
  systematically silence low-expressiveness parties (e.g., a patient with
  reduced capacity), violating the very justice principle it claims to
  protect. This poses higher demands for \emph{mechanism-level fairness}.
\end{enumerate}

% P5 方法预告
To handle the issues above, inspired by the theory of generalized Nash
equilibrium with shared constraints \cite{facchinei2010,shared2017}, we
introduce \textbf{EthicalGuard}, a hybrid LLM + numerical-solver negotiation
framework, with three core properties considered: \emph{provable floors},
\emph{justified solution concept}, and \emph{fair reliability}. Specifically:
\begin{itemize}
  \item[(i)] we formulate ethics as a regularized GNE problem and prove a
  floor-guarantee theorem from the KKT conditions (Theorem~\ref{thm:floor}),
  with a designed floor-activation experiment showing the constraint binds;
  \item[(ii)] we argue the modeling choice of GNE over NBS/Stackelberg/
  cooperative games (Section~3.3), and declare the numerical approximation
  (entropy regularization, $\eta>0$; balance penalty $\gamma$) precisely;
  \item[(iii)] we introduce a reliability-weight floor
  $w_i\ge\varepsilon$ and decouple information reliability from stance
  representativeness, preventing self-reinforcing silencing of weak parties.
\end{itemize}

% P6 贡献三件套
In summary, our main contributions are as follows:
\begin{itemize}
  \item We investigate the problem of provable ethical floors in multi-agent
  medical decision-making and point out that static normative filters lack
  equilibrium-level guarantees, that solution-concept choice is undefended,
  and that reliability weighting can violate justice.
  \item We propose \textbf{EthicalGuard}, which realizes ``the ethical floor
  is a shared constraint'' as a regularized GNE with full KKT validation,
  hard veto, conflict reporting, and a catfish agent; we prove the
  floor-guarantee theorem, justify the GNE modeling choice, and introduce the
  reliability-weight floor for fairness.
  \item Extensive experiments on five medical-ethics datasets
  (MedEthicEval, MedEthicsQA, PrinciplismQA, VITAL, LLMEval-Med) reveal
  convergence with KKT residuals $\sim10^{-7}$, component contributions via
  ablations, resilience under stress, and the designed validation of
  floor activation, collusion resistance, and fairness.
\end{itemize}

% ---------- 2. Related Work ----------
\section{Related Work}
\subsection{Medical LLM Ethics Reasoning}
Uncertainty-aware models quantify diagnostic uncertainty and abstention
\cite{confidx,entropy2026}; clinical selective prediction argues for actions
beyond binary abstain \cite{limits2025}. These address \emph{epistemic}
soundness of a single model, not the reconciliation of \emph{conflicting
values} among stakeholders. Ethical reasoning in medical LLMs is benchmarked
by MedEthicEval (violation/equilibrium/priority dilemmas), MedEthicsQA,
PrinciplismQA, and VITAL (value distribution and steerability)
\cite{medethiceval,medethicsqa,principlismqa,vital}. However, these works
focus on measuring or calibrating a single model's ethics, overlooking the
\emph{multi-party negotiation} structure of real ethical dilemmas.

\subsection{Multi-Agent Medical Decision-Making}
Multi-agent systems structure clinical reasoning: Aegle \cite{aegle} models
SOAP state with observation masks; EmoMAS \cite{emomas} orchestrates agents
with Bayesian reliability weights; MedLA \cite{medla} chains logical trees;
KAMAC \cite{kamac} recruits agents dynamically; Graph-of-States
\cite{graphofstates} drives FSM transitions. These provide building blocks
for multi-perspective reasoning, but none formalizes a game-theoretic
equilibrium over conflicting values, nor a non-negotiable floor.
\emph{Many LLMs Are More Utilitarian Than One} \cite{utilitarian2025} shows
aggregation shifts group decisions toward utilitarianism---yet voting neither
models strategic stakeholders nor enforces shared constraints. Our work
bridges this gap by adding a GNE layer (shared constraints = ethical floors)
with full KKT validation and auditable syllogistic reasoning.

\subsection{Normative Constraints and Game-Theoretic Ethics}
Constitution-aligned ``superego'' layers \cite{superego2025}, governance
frameworks \cite{governance2025}, and plug-in conscience mechanisms share the
intuition that some constraints are non-negotiable, but apply them
statically (at the agent or system level). Game-theoretic treatments of
ethics exist at the abstract level: Nash bargaining and fair-division games
model conflict resolution; GNE extends Nash to shared coupling constraints
\cite{facchinei2010,shared2017}. A recent review of ethical issues in
multi-agent healthcare AI \cite{frontiers2026} highlights compound opacity,
motivating our auditable syllogism-grounded negotiation. Our work instantiates
the philosophical claim ``floors are non-negotiable'' as GNE shared
constraints in a medical LLM system, and handles the bad-equilibrium problem
via veto, conflict reports, and a catfish agent.

% ---------- 3. Methodology ----------
\section{Methodology}
The overview of EthicalGuard is detailed in Fig.~\ref{fig:arch}: a five-module
pipeline (M0 conflict detection $\to$ M1 dynamic assembly $\to$ M2 negotiation
loop $\to$ M3 arbitration $\to$ M4 GNE refinement $\to$ M5 resilience), with
LLM-generated (or rule-based) stances feeding a numerical GNE solver that
refines principle weights under full KKT validation; every proposal is emitted
as a syllogism for auditability.
\begin{figure}[t]
  \centering
  % \includegraphics[width=\columnwidth]{architecture.png}
  \caption{EthicalGuard architecture: M0 gates negotiation via an ERS risk
  signal; M1 assembles the five stakeholders plus a catfish agent; M2 runs the
  state machine with Bayesian orchestration; M3 enforces L1 hard vetoes and
  CAMP arbitration; M4 solves the regularized GNE in weight space with full
  KKT validation; M5 stresses the outcome and verifies recovery.}
  \label{fig:arch}
\end{figure}

\subsection{Problem Formulation: Ethics as Regularized GNE}
\label{sec:formulation}
\subsubsection{Game structure}
Let $\mathcal{I} = \{1,\dots,N\}$ be the stakeholder agents (physician,
ethics committee, patient, family, hospital administration) and let the
catfish agent $\mathcal{C}$ be a non-strategic dissenter that does not enter
the collective vector (see Section~\ref{sec:catfish}). Each strategic agent
$i$ holds strategy $x_i = (t_i, \alpha_i)$: treatment intensity
$t_i \in \{0,1,2,3\}$ (carrying resource usage $\rho(t_i)$ with
$0/1/2/3 \to 0/0.3/0.6/1.0$) and principle weights
$\alpha_i \in \Delta^4$ over (beneficence, non-maleficence, autonomy,
justice). The individual objective is
\begin{equation}
  J_i(x_i) = \alpha_i \cdot s_i - \mathrm{soft\_cost}_i \cdot \alpha_i
  + \gamma\,\mathrm{FDBI}(v) - \eta \sum_{k} \alpha_{ik}\log\alpha_{ik},
  \label{eq:Ji}
\end{equation}
where $s_i$ encodes the agent's stance satisfaction, $\mathrm{soft\_cost}_i$
its soft constraints, $v = \sum_i w_i \alpha_i / \sum_i w_i$ the collective
vector with $w_i$ the orchestration weight, $\gamma\ge0$ the balance-penalty
weight, and $\eta>0$ the entropy-regularization strength. The
\textbf{shared coupling constraints} are
\begin{equation}
  g_{res}(x) = \sum_{i} \rho(t_i) - cap \le 0,
  \qquad
  g_{floor,k}(v) = floor_k - v_k \le 0, \quad k \in \{B,N,A,J\},
  \label{eq:constraints}
\end{equation}
solved via a two-timescale primal--dual method with full KKT validation
(stationarity, primal/dual feasibility, complementary slackness).

\subsubsection{Exact solution concept: entropy-regularized GNE}
\label{sec:solution-concept}
We emphasize the precise solution concept: Eq.~(\ref{eq:Ji}) with $\eta>0$
and $\gamma>0$ defines an \emph{entropy-regularized} generalized Nash game,
whose solution is the unique interior fixed point of the damped best-response
dynamics (the entropy term makes $J_i$ strongly concave in $\alpha_i$ given
the collective $v$). As $\eta \to 0$ and $\gamma \to 0$, the solution tends
to the classical GNE of the limiting game. We therefore report KKT residuals
of the regularized problem and state all guarantees (Theorem~\ref{thm:floor})
for the regularized solution, which is the object our system actually
computes. This declaration preempts the objection that ``the implemented
game is not the textbook GNE'': it is a well-defined regularization of it,
with an explicit limit correspondence.

\subsection{Floor-Guarantee Theorem}
\label{sec:theorem}
The semantic core of our formulation is that a \emph{generalized} Nash
equilibrium---unlike Nash---requires all agents to jointly satisfy the
coupling constraints $g(x)\le0$. We interpret $g_{res}$ as the resource
reality of the joint treatment plan and $g_{floor,k}$ as the four-principle
floors evaluated on the collective vector. \textbf{Negotiation is free within
the floors; the floors are non-negotiable.} This claim is not asserted but
proved:

\begin{theorem}[Floor guarantee]
\label{thm:floor}
Let $v^{*}$ be the collective vector returned by the primal--dual solver of
Section~\ref{sec:formulation}, and let $\mathrm{KKT}_{total}$ be its total
KKT residual (the sum of stationarity, primal, dual, and complementary
slackness residuals). Then for every principle $k$,
\begin{equation}
  v^{*}_k \;\ge\; floor_k \;-\; \mathrm{KKT}_{total}.
  \label{eq:floor-guarantee}
\end{equation}
Moreover, if the complementary-slackness residual for the $k$-th floor
vanishes ($\lambda_k\, g_{floor,k}(v^{*}) \to 0$) with multiplier
$\lambda_k > 0$, then the constraint binds at equality up to tolerance:
$v^{*}_k = floor_k + O(\mathrm{KKT}_{total})$.
\end{theorem}

\begin{proof}
The total KKT residual is the sum of four non-negative terms, one of which is
the primal feasibility residual
$\mathrm{primal} = \max(0, g_{res}) + \sum_k \max(0, g_{floor,k}(v^{*}))$.
Since every term is non-negative, $\mathrm{primal} \le \mathrm{KKT}_{total}$.
Hence for each $k$,
$\max(0, floor_k - v^{*}_k) \le \mathrm{KKT}_{total}$, which is equivalent
to $v^{*}_k \ge floor_k - \mathrm{KKT}_{total}$, proving the first claim. For
the second claim, the complementary-slackness residual for floor $k$ is
$|\lambda_k\, g_{floor,k}(v^{*})| \le \mathrm{KKT}_{total}$. If
$\lambda_k > 0$ (the constraint is active) and the residual vanishes, then
$g_{floor,k}(v^{*}) \to 0$, i.e., $v^{*}_k \to floor_k$. ∎
\end{proof}

\begin{remark}[Non-triviality and activation]
\label{rem:activation}
Theorem~\ref{thm:floor} guarantees the floor in every converged run, but this
guarantee is only \emph{informative} if the constraint is \emph{active}: if
the default floors (0.15--0.20) are far below the balance target (0.25), the
balance penalty alone may keep $v$ above the floors, making the constraint
never bind. We therefore design a floor-activation experiment (RQ6,
Section~\ref{sec:rq6}) with extreme-stance scenarios that push a principle
below its floor, and verify that (i) the multiplier $\lambda_k$ becomes
positive, (ii) the solution is pulled back to $v^{*}_k = floor_k$, and (iii)
the veto mechanism triggers on direct violations. Without activation
evidence, Theorem~\ref{thm:floor} is vacuous in practice.
\end{remark}

\begin{proposition}[Violation exclusion]
\label{prop:veto}
Under the arbitrator's L1 hard veto (Section~\ref{sec:arbitration}), any
strategy profile violating a floor constraint is not a feasible equilibrium
output: either the solver keeps it infeasible (primal residual $>0$) or the
arbitrator excludes or downgrades it. Hence no floor-violating outcome
survives the constraint--enforcement loop.
\end{proposition}
\begin{proof}
By construction of the enforcement loop: the GNE solver certifies feasibility
and reports shadow prices; the arbitrator's L1 veto (which excludes
resource-type constraints, handled by the resource-cap correction) rejects any
proposal with a positive floor violation; the catfish agent injects structured
dissent against the most-neglected principle. A floor-violating profile is
therefore either not an equilibrium of the feasible game or is vetoed before
becoming the output. ∎
\end{proof}

\subsection{Modeling Choice: Why GNE?}
\label{sec:modeling-choice}
A critic may ask: why GNE rather than a solution concept with axiomatic
foundations? We answer by exclusion:

\begin{itemize}
  \item \textbf{Nash bargaining (NBS)}: NBS is a \emph{cooperative} solution
  concept that assumes a binding agreement and a known feasible set. In
  clinical ethics there is no binding contract, and NBS axioms are
  contestable in this domain (independence of irrelevant alternatives fails
  when adding an option should change the ethical judgment). Crucially, NBS
  cannot express a \emph{shared, non-negotiable} floor: the bargaining
  solution selects a Pareto point inside the feasible set, treating
  constraints as exogenous filters that can themselves be bargained over. A
  floor is not something to bargain about.
  \item \textbf{Stackelberg}: presupposes a leader--follower hierarchy; but
  in ethics, \emph{who leads} (physician authority vs.\ patient autonomy) is
  itself the contested question. Stackelberg turns authority into an input
  rather than a negotiated outcome.
  \item \textbf{Cooperative games (core/Shapley)}: assume (transferable)
  utility, which principle weights violate---one cannot transfer autonomy
  weight to justice. For non-transferable utility, the core may be empty.
  \item \textbf{GNE (our choice)}: non-cooperative, with \emph{shared
  coupling constraints}. This matches the ethical semantics exactly:
  negotiation is free within the floors (each agent optimizes its own
  objective), and the floors are part of the equilibrium structure itself
  (not an exogenous filter). The KKT multipliers $\lambda_{floor,k}$ carry
  ethical meaning as shadow prices: they quantify the marginal value of
  relaxing a floor, i.e., which principle is being ``pulled'' by the
  negotiation.
\end{itemize}
We further declare the numerical approximation explicitly
(Section~\ref{sec:solution-concept}): the solver is an entropy-regularized
GNE, and all theorems and metrics refer to this regularized object.

\subsection{System Architecture}
\label{sec:arch}
EthicalGuard runs a five-module pipeline: \textbf{(M0) conflict detection}
gates negotiation via an ERS risk signal and a supervised violation detector;
\textbf{(M1) dynamic assembly} recruits the five stakeholders plus the catfish
agent; \textbf{(M2) negotiation loop} runs a state machine (S0 fact anchoring
$\to$ S1 assembly $\to$ S2 propose $\to$ S3 alignment $\to$ S4 stress $\to$
S5 arbitration $\to$ S6 recovery, with R backtracking on drift) and updates
reliability weights via Bayesian orchestration; \textbf{(M3) arbitration}
\label{sec:arbitration} enforces L1 hard vetoes (excluding resource-type
constraints, handled by the resource-cap correction) and CAMP
(KEEP/REFUSE/NEUTRAL) evidence arbitration; \textbf{(M4) GNE refinement}
solves the regularized equilibrium in weight space with full KKT validation;
\textbf{(M5) resilience} stresses the negotiated outcome across an intensity
grid and verifies recovery.

\subsection{Fair Reliability: Weight Floors and Decoupling}
\label{sec:fairness}
\begin{definition}[Reliability weight and stance weight]
Let $r_i$ be the Bayesian reliability of agent $i$ (Beta posterior mean,
updated by whether its proposal is adopted), used \emph{only} to trigger
arbitration and to report information quality. Let $w_i$ be the stance weight
used to form the collective vector $v = \sum_i w_i \alpha_i / \sum_i w_i$. We
set
\begin{equation}
  w_i \;=\; \frac{\max(r_i,\ \varepsilon)}{\sum_j \max(r_j,\ \varepsilon)},
  \qquad \varepsilon \ge \frac{\beta_0}{N},
  \label{eq:weight-floor}
\end{equation}
for a minimum weight $\varepsilon$ (with $\beta_0$ a small constant, e.g.,
0.5), and we \emph{decouple} $r_i$ from $w_i$: reliability reflects
information quality, not moral weight.
\end{definition}

\begin{proposition}[No silencing]
\label{prop:silencing}
With the weight floor of Eq.~(\ref{eq:weight-floor}), no agent's stance
weight can vanish: $w_i \ge \varepsilon / N > 0$ (since
$\max(r_j,\varepsilon) \le 1$ for every $j$, the denominator is at most $N$).
Consequently, the collective vector and the floor constraint
$g_{floor,k}(v)$ always retain a non-vanishing contribution from every
stakeholder, preventing the self-reinforcing marginalization loop in which a
low-expressiveness agent is progressively silenced and then excluded from
floor determination.
\end{proposition}
\begin{proof}
Direct from Eq.~(\ref{eq:weight-floor}): the numerator is at least
$\varepsilon$ for every $i$, and the denominator
$\sum_j \max(r_j,\varepsilon) \le N$ (each term is at most 1), so
$w_i \ge \varepsilon / N > 0$. The floor constraint is evaluated on
$v=\sum_i w_i\alpha_i/\sum_i w_i$, a convex combination with all coefficients
bounded below by $\varepsilon/N$; hence every agent's stance has weight at
least $\varepsilon/N$ in the floor determination, and in particular no agent
can be silenced as $r_i \to 0$. ∎
\end{proof}

\begin{remark}
The decoupling has a normative justification: information reliability
(how well an agent articulates evidence) is not moral weight (how much its
values count). A patient with reduced capacity may articulate poorly yet her
autonomy floor is exactly what the system must protect; weighting her by
reliability alone would invert the protective intent of the floor. We
additionally report the trajectory of $w_{patient}$ under low-capacity
scenarios (RQ8) to empirically confirm the absence of collapse.
\end{remark}

\subsection{Catfish Dissent}
\label{sec:catfish}
The catfish agent $\mathcal{C}$ is a non-strategic dissenter: it does not
occupy resources, does not enter the collective vector, and is excluded from
GNE solving (its role is observational dissent). It injects structured
objections targeting the currently most-neglected principle
($\arg\min_k v_k$). We treat catfish as a \emph{heuristic component of the
constraint--enforcement loop} (Proposition~\ref{prop:veto}) rather than a
game-theoretic perturbation, and validate its collusion-breaking effect
empirically in RQ7 (Section~\ref{sec:rq7}). This honest positioning avoids
claiming a formal anti-collusion theorem we do not have.

\subsection{Hybrid LLM + Solver Stances}
Each agent's initial proposal (treatment level and principle weights) is
generated by an LLM (or by deterministic rules in the reproducible baseline
mode) from the scenario's anchored facts; the GNE solver then refines the
principle weights in $\Delta^4$, coupling the discrete scheme space via the
resource mapping $\rho(t)$. This hybrid avoids LLM imprecision in numerical
equilibrium computation while retaining LLM semantic richness in stance
formation. All rationales are emitted as syllogisms (major premise, minor
premise, conclusion), making every proposal auditable at the format level;
we explicitly do not claim audited logical correctness of LLM syllogisms
(see Limitations).

% ---------- 4. Experiments ----------
\section{Experiments}
We design experiments to answer the following questions:
\begin{itemize}
  \item \textbf{RQ1} (convergence): Does the regularized GNE negotiation
  converge to a stable, floor-respecting consensus?
  \item \textbf{RQ2} (components): What are the effects of each component
  (GNE refinement, catfish, arbitration, balance regularization)?
  \item \textbf{RQ3} (sensitivity): How sensitive is the outcome to the
  balance penalty $\gamma$ and to reward/harm weights?
  \item \textbf{RQ4} (resilience): Does the negotiated outcome recover under
  stress, and at what intensity are principles abandoned?
  \item \textbf{RQ5} (modes): How do deterministic-rule and LLM-stance modes
  differ in convergence and floor enforcement?
  \item \textbf{RQ6} (floor activation): Are the floor constraints
  \emph{active} (binding) in extreme-stance scenarios, as required for
  Theorem~\ref{thm:floor} to be non-vacuous?
  \item \textbf{RQ7} (bad equilibria): Do catfish and floor enforcement
  break collusion in adversarially constructed scenarios?
  \item \textbf{RQ8} (fairness): Does the reliability-weight floor prevent
  the silencing of low-expressiveness stakeholders?
\end{itemize}

\subsection{Experimental settings}
\subsubsection{Datasets}
We evaluate on five datasets: \textbf{MedEthicEval} (violation/equilibrium/
priority dilemmas), \textbf{MedEthicsQA} (MCQ + open-ended), \textbf{PrinciplismQA}
(principle-grounded QA), \textbf{VITAL} (value distribution and steerability),
and \textbf{LLMEval-Med} (medical LLM evaluation). The first four provide
scenarios for negotiation; LLMEval-Med is used for stance-quality checks.
Statistics are summarized in Table~\ref{tab:data}.
\begin{table}[t]
  \centering
  \caption{Statistics of the five datasets (scenario counts after
  preparation).}
  \label{tab:data}
  \begin{tabular}{lcc}
    \toprule
    Dataset & Scenarios & Task type \\
    \midrule
    MedEthicEval & TBD & violation / equilibrium / priority \\
    MedEthicsQA & TBD & MCQ + open-ended \\
    PrinciplismQA & TBD & principle-grounded QA \\
    VITAL & TBD & value distribution / steerability \\
    LLMEval-Med & TBD & medical LLM evaluation \\
    \bottomrule
  \end{tabular}
\end{table}
\subsubsection{Baselines}
We compare against two families: (i) \emph{single-model} baselines---standard
LLM diagnosis, entropy-threshold selective prediction \cite{entropy2026},
uncertainty-aware diagnosis \cite{confidx}; and (ii) \emph{multi-agent}
baselines---majority/plurality voting over LLM stances
\cite{utilitarian2025}, GraphDx-style cost-sensitive reasoning \cite{graphdx},
and MAI-DxO-style orchestrator \cite{maidxo}. We further ablate EthicalGuard
itself (RQ2).
\subsubsection{Implementation details}
Negotiation runs under two modes: deterministic rules (reproducible baseline,
zero external API) and LLM stances (local model, temperature 0.0 for
reproducibility). Metrics: GNE KKT residual (per-scenario distribution),
stakeholder satisfaction (cosine between own weights and collective vector),
residual conflict (mean pairwise L1 distance), ERS conflict scores, floor
violation counts, and floor-activation multipliers.

\subsection{Convergence (RQ1)}
\label{sec:rq1}
We observe that: 1) on real negotiation trajectories, the regularized GNE
solver converges with per-scenario KKT residuals on the order of $10^{-7}$
(e.g., $1.69\times10^{-7}$, $1.16\times10^{-7}$ in representative
trajectories), well below the tolerance $10^{-3}$; 2) mean stakeholder
satisfaction lies in 0.82--0.99 across scenarios, with patient satisfaction
the most variable (0.82 in resource-pressured cases); 3) residual conflict
decreases across rounds, indicating collective convergence; 4) by
Theorem~\ref{thm:floor}, every converged output satisfies the floors up to
numerical tolerance---floor violations are zero in standard scenarios
(Table~\ref{tab:main}).
\begin{table}[t]
  \centering
  \caption{Convergence and satisfaction (representative trajectories; full
  distributions reported in appendix).}
  \label{tab:main}
  \begin{tabular}{lccc}
    \toprule
    Scenario & KKT residual & Mean satisfaction & Floor violations \\
    \midrule
    PQ-1-2 & $1.69\times10^{-7}$ & 0.95 & 0 \\
    PQ-3-7 & $1.16\times10^{-7}$ & 0.88 & 0 \\
    (full set) & TBD & TBD & 0 \\
    \bottomrule
  \end{tabular}
\end{table}

\subsection{Component ablations (RQ2)}
\label{sec:rq2}
We ablate four components: GNE refinement (w/o GNE: naive Bayesian collective),
catfish (w/o Catfish), arbitration (w/o arbitration), and balance
regularization (w/o balance). Table~\ref{tab:ablation} reports FDBI, PCI,
KKT rate, arbitration rate, feasibility, floor violations, and minimum
satisfaction. We observe that: 1) removing GNE refinement raises floor
violations (0.125), showing the solver contributes to floor enforcement;
2) removing arbitration collapses the KKT success rate (30\%) and feasibility
(30\%), confirming the arbitrator's role in convergence; 3) removing the
balance regularizer slightly reduces FDBI; 4) removing catfish changes the
arbitration rate, consistent with its dissent role.
\begin{table}[t]
  \centering
  \caption{Ablation results (mode=local).}
  \label{tab:ablation}
  \begin{tabular}{lcccccc}
    \toprule
    Variant & FDBI & PCI & KKT rate & Arbitration & Feasibility & Floor viol. \\
    \midrule
    MANE (full) & 0.882 & 0.075 & 100\% & 70\% & 100\% & 0.000 \\
    w/o GNE & 0.800 & 0.127 & nan\% & 70\% & 100\% & 0.125 \\
    w/o Catfish & 0.899 & 0.063 & 100\% & 30\% & 100\% & 0.000 \\
    w/o Arbitration & 0.912 & 0.055 & 30\% & 0\% & 30\% & 0.000 \\
    w/o Balance & 0.880 & 0.076 & 100\% & 70\% & 100\% & 0.000 \\
    \bottomrule
  \end{tabular}
\end{table}

\subsection{Sensitivity (RQ3)}
\label{sec:rq3}
We sweep the balance penalty $\gamma$ (0.0--1.0). We observe that FDBI
increases monotonically with $\gamma$ (0.864 $\to$ 0.875), while floor
violations remain zero at all settings---the balance regularizer improves
collective balance without breaking floors (Table~\ref{tab:gamma}).
\begin{table}[t]
  \centering
  \caption{$\gamma$ sensitivity (mode=local, 5 scenarios).}
  \label{tab:gamma}
  \begin{tabular}{lcccc}
    \toprule
    $\gamma$ & FDBI & PCI & Floor viol. & Min satisfaction \\
    \midrule
    0.00 & 0.864 & 0.084 & 0.000 & 0.807 \\
    0.20 & 0.866 & 0.082 & 0.000 & 0.807 \\
    0.40 & 0.869 & 0.081 & 0.000 & 0.807 \\
    0.60 & 0.871 & 0.079 & 0.000 & 0.807 \\
    0.80 & 0.873 & 0.078 & 0.000 & 0.807 \\
    1.00 & 0.875 & 0.077 & 0.000 & 0.807 \\
    \bottomrule
  \end{tabular}
\end{table}

\subsection{Resilience (RQ4)}
\label{sec:rq4}
We stress the negotiated outcome with seven perturbation types across an
intensity grid $\{0.25,0.5,0.75,1.0\}$ and measure three-level resilience
(consistency, robustness, recoverability), abandonment events, and the
critical abandonment intensity $\tau$ per principle. In LLM-stance mode, we
observe real abandonment evidence---e.g., in scenario PQ-1-2,
$abandoned=\{justice: authority,\ autonomy: disease\}$ with
$\tau_{justice}=0.25$ and $\tau_{autonomy}=0.5$ under authority/disease
stress---whereas the rule mode serves as a consistency baseline
($baseline\_only=true$) with near-perfect consistency and no genuine
abandonment. We note that $l_2$-robustness values are low (0.05--0.10) while
the overall resilience is high (0.77--0.82); we interpret this as the
consistency/recoverability terms dominating, and we report the metric
decomposition transparently (Table~\ref{tab:resilience}).
\begin{table}[t]
  \centering
  \caption{Resilience under stress (mode=local).}
  \label{tab:resilience}
  \begin{tabular}{lcccc}
    \toprule
    Scenario & Consistency & Robustness & Recoverability & Overall \\
    \midrule
    PQ-1-2 & 0.964 & 0.102 & 0.566 & 0.822 \\
    PQ-3-7 & 0.975 & 0.053 & 0.321 & 0.769 \\
    (full set) & TBD & TBD & TBD & TBD \\
    \bottomrule
  \end{tabular}
\end{table}

\subsection{Mode comparison (RQ5)}
\label{sec:rq5}
We compare deterministic-rule and LLM-stance modes on convergence, floor
enforcement, and resilience. Rule mode is fully reproducible and shows
consistency-baseline behavior; LLM mode exhibits genuine concession
(abandonment events under stress) while retaining convergence quality.
Full pairwise comparison is reported in Table~\ref{tab:modes}.
\begin{table}[t]
  \centering
  \caption{Rule vs.\ LLM modes.}
  \label{tab:modes}
  \begin{tabular}{lccc}
    \toprule
    Mode & KKT residual & Floor viol. & Abandonment evidence \\
    \midrule
    rule & TBD & 0 & none (baseline) \\
    LLM (local) & $\sim10^{-7}$ & 0 & yes (RQ4) \\
    \bottomrule
  \end{tabular}
\end{table}

\subsection{Floor activation (RQ6)}
\label{sec:rq6}
To show Theorem~\ref{thm:floor} is non-vacuous, we construct extreme-stance
scenarios: a party insists on a principle weight below its floor (e.g., a
proxy pushing non-maleficence to 0.10 when $floor_N=0.20$). We verify
(i) the floor multiplier $\lambda_{floor,k}$ becomes positive, (ii) the
collective vector is pulled back to $v^{*}_k = floor_k$, and (iii) direct
violations trigger the L1 veto. Table~\ref{tab:floor} reports activation
statistics.
\begin{table}[t]
  \centering
  \caption{Floor activation in extreme-stance scenarios.}
  \label{tab:floor}
  \begin{tabular}{lcccc}
    \toprule
    Scenario & $\lambda_{k}>0$ & $v^{*}_k$ & Veto triggered & Bound \\
    \midrule
    (designed) & TBD & TBD & TBD & TBD \\
    \bottomrule
  \end{tabular}
\end{table}

\subsection{Bad-equilibrium stress (RQ7)}
\label{sec:rq7}
We construct collusion scenarios: four agents coordinate to suppress the
patient's autonomy floor (e.g., all proposing autonomy $=0.10$). We verify
that the catfish agent breaks the silent consensus and that the floor
constraint/veto prevents a floor-violating outcome from surviving
(Proposition~\ref{prop:veto}). We report whether the colluding profile is
excluded or destabilized.
\begin{table}[t]
  \centering
  \caption{Collusion scenarios.}
  \label{tab:collusion}
  \begin{tabular}{lccc}
    \toprule
    Scenario & Catfish triggered & Floor violated & Survives? \\
    \midrule
    (designed) & TBD & TBD & TBD \\
    \bottomrule
  \end{tabular}
\end{table}

\subsection{Fairness (RQ8)}
\label{sec:rq8}
We run low-capacity patient scenarios (reduced clarity) and report (i) the
trajectory of $w_{patient}$ across negotiation rounds with and without the
weight floor of Eq.~(\ref{eq:weight-floor}); (ii) the divergence between the
collective vector and the patient's stance; (iii) whether the patient's
autonomy floor is protected. We expect the weight floor to keep
$w_{patient} \ge \varepsilon$ and prevent the silencing collapse
(Proposition~\ref{prop:silencing}).
\begin{table}[t]
  \centering
  \caption{Fairness under low-expressiveness patients.}
  \label{tab:fairness}
  \begin{tabular}{lccc}
    \toprule
    Setting & $w_{patient}$ (final) & $v$ vs.\ patient divergence & Floor protected \\
    \midrule
    w/o weight floor & TBD & TBD & TBD \\
    w/ weight floor & TBD & TBD & TBD \\
    \bottomrule
  \end{tabular}
\end{table}

\subsection{Case study}
Fig.~\ref{fig:case} shows one representative negotiation: the detained trauma
patient BB with a correctional officer observing. Five stakeholders propose
treatment intensities and principle weights; the ethics committee invokes the
L1 veto on resource grounds; the GNE solver refines the collective vector;
the conflict report recommends a 24-hour review; satisfaction converges above
0.8. This trajectory illustrates auditable, floor-respecting,
multi-stakeholder consensus.
\begin{figure}[t]
  \centering
  % \includegraphics[width=\columnwidth]{case.png}
  \caption{Representative negotiation trajectory for the BB scenario.}
  \label{fig:case}
\end{figure}

% ---------- 5. Conclusion ----------
\section{Conclusion}
We formulated medical ethics decision-making as a regularized generalized
Nash equilibrium problem in which the ethical floor is realized as shared
constraints, proved a floor-guarantee theorem from the KKT conditions, and
justified the GNE modeling choice by exclusion of NBS/Stackelberg/cooperative
games. We introduced a reliability-weight floor decoupling information
reliability from stance representativeness, preventing the silencing of
low-expressiveness stakeholders. On five medical-ethics datasets, EthicalGuard
converges with KKT residuals $\sim10^{-7}$, satisfies floors, ablated
components confirm contributions, and stress tests document recovery and
principle abandonment.
\paragraph{Limitations.}
(i) The floor guarantee is for the entropy-regularized GNE, not the limiting
classical GNE; (ii) catfish collusion-breaking is validated empirically, not
by an anti-collusion theorem; (iii) ``auditable syllogisms'' guarantee format
auditability, not logical correctness of LLM-generated premises; (iv) the
fairness decoupling rests on a normative claim (reliability $\neq$ moral
weight) that we state explicitly and validate mechanically (weight-floor
effect) but not clinically; (v) the weight floor $\varepsilon$ and balance
penalty $\gamma$ are design parameters requiring sensitivity analysis.
\paragraph{Future work.}
We plan to validate the floor guarantees against clinical-ethics expert
judgments (gold-standard alignment), scale to real patient narratives, and
extend the fairness analysis to multi-round delegation and proxy-decision
settings.

\section*{CRediT authorship contribution statement}
<作者贡献>。
\section*{Declaration of competing interest}
<利益声明>。
\section*{Data availability}
The four medical-ethics datasets (MedEthicEval, MedEthicsQA, PrinciplismQA,
VITAL) and LLMEval-Med are publicly available; code and reproducible
pipelines are released with the paper.

\bibliographystyle{plain}
\bibliography{references}

\end{document}
